% GENERATED FILE. Edit canonical lessons, then run npm run generate:books.
% canonical-source-hash: fnv1a64:57b5365f0a6f5cf7
% canonical-lessons: HI-C239-camakdar, HI-C239-dhudhla, HI-C239-sunsan, HI-C239-bhirbhar, HI-C239-abad

\chapter{Shiny, Hazy, Deserted, Crowded, Inhabited}
\label{ch:hi-shiny-hazy-deserted-crowded-inhabited}
\hlchaptermodality{239}

\begin{chapteropening}
\textbf{By the end of this chapter:} \emph{I can say shiny, hazy, deserted, crowded and inhabited.}

Complete the last lesson of chapter 239: I can say shiny, hazy, deserted, crowded and inhabited.
\end{chapteropening}

\section[\texorpdfstring{\hi{चमकदार} (camakdār)}{चमकदार (camakdār)}]{\hi{चमकदार} (camakdār) — shiny}

\label{lesson:HI-C239-camakdar}



\begin{quote}
Before the new one: say the Hindi for sticky, then the Hindi for rough.
\end{quote}

\subsection*{You'll want to know: \hi{चमकदार}}
\textbf{\hi{चमकदार}} — \emph{camakdār} — \textquotedblleft{}shiny\textquotedblright{}.

\begin{cousinweb}
\textbf{\hi{चमकदार}} (\emph{camakdār}) is \emph{camak}, \textquotedblleft{}shine\textquotedblright{}, from \emph{camaknā}, \textquotedblleft{}to glitter\textquotedblright{}, which CDIAL lists under a reconstructed \emph{*camakka}, \textquotedblleft{}sudden movement\textquotedblright{} (CDIAL 4676), with the Persian ending \emph{-dār}, \textquotedblleft{}holder\textquotedblright{}.
\end{cousinweb}

\begin{grammarlens}[title={What you've built}]
One more word for this chapter.
\end{grammarlens}

\subsection*{Your turn}
\begin{itemize}
\raggedright
  \item \emph{Say it:} \emph{camakdār}
  \item \emph{Say it:} \emph{camakdār}, once more
  \item \emph{Say it:} say \emph{camakdār}
  \item \emph{Recall:} say the Hindi for sticky, then the Hindi for rough, then say \emph{camakdār} again
\end{itemize}

\begin{tcolorbox}[breakable,colback=teal!4,colframe=teal!35!black,title={Before you move on}]
What does \hi{चमकदार} mean? (\textquotedblleft{}Shiny\textquotedblright{}.) Say it once more.
\end{tcolorbox}
\section[\texorpdfstring{\hi{धुंधला} (dhũdhlā)}{धुंधला (dhũdhlā)}]{\hi{धुंधला} (dhũdhlā) — hazy, blurred}

\label{lesson:HI-C239-dhudhla}



\begin{quote}
Before the new one: say the Hindi for rough, then the Hindi for shiny.
\end{quote}

\subsection*{You'll want to know: \hi{धुंधला}}
\textbf{\hi{धुंधला}} — \emph{dhũdhlā} — \textquotedblleft{}hazy, blurred\textquotedblright{}.

\begin{cousinweb}
\textbf{\hi{धुंधला}} (\emph{dhũdhlā}) is a tadbhava, a word inherited and reshaped by sound change. It comes from \emph{*dhūmāndha}, \textquotedblleft{}dim-sighted\textquotedblright{}, a form CDIAL reconstructs from its descendants (CDIAL 6858). It is built from \emph{dhūma}, \textquotedblleft{}smoke\textquotedblright{}, and \emph{andha}, \textquotedblleft{}blind\textquotedblright{}.
\end{cousinweb}

\begin{grammarlens}[title={What you've built}]
One more word for this chapter.
\end{grammarlens}

\subsection*{Your turn}
\begin{itemize}
\raggedright
  \item \emph{Say it:} \emph{dhũdhlā}
  \item \emph{Say it:} \emph{dhũdhlā}, once more
  \item \emph{Say it:} say \emph{dhũdhlā}
  \item \emph{Recall:} say the Hindi for rough, then the Hindi for shiny, then say \emph{dhũdhlā} again
\end{itemize}

\begin{tcolorbox}[breakable,colback=teal!4,colframe=teal!35!black,title={Before you move on}]
What does \hi{धुंधला} mean? (\textquotedblleft{}Hazy, blurred\textquotedblright{}.) Say it once more.
\end{tcolorbox}
\section[\texorpdfstring{\hi{सुनसान} (sunsān)}{सुनसान (sunsān)}]{\hi{सुनसान} (sunsān) — deserted}

\label{lesson:HI-C239-sunsan}



\begin{quote}
Before the new one: say the Hindi for shiny, then the Hindi for hazy.
\end{quote}

\subsection*{You'll want to know: \hi{सुनसान}}
\textbf{\hi{सुनसान}} — \emph{sunsān} — \textquotedblleft{}deserted\textquotedblright{}.

\begin{cousinweb}
The origin of \textbf{\hi{सुनसान}} (\emph{sunsān}) is uncertain. It may be built on \emph{sūn}, \textquotedblleft{}empty\textquotedblright{}, which CDIAL lists under Sanskrit \emph{śūnya}, \textquotedblleft{}empty\textquotedblright{} (CDIAL 12567), but CDIAL does not list \emph{sunsān} itself.
\end{cousinweb}

\begin{grammarlens}[title={What you've built}]
One more word for this chapter.
\end{grammarlens}

\subsection*{Your turn}
\begin{itemize}
\raggedright
  \item \emph{Say it:} \emph{sunsān}
  \item \emph{Say it:} \emph{sunsān}, once more
  \item \emph{Say it:} say \emph{sunsān}
  \item \emph{Recall:} say the Hindi for shiny, then the Hindi for hazy, then say \emph{sunsān} again
\end{itemize}

\begin{tcolorbox}[breakable,colback=teal!4,colframe=teal!35!black,title={Before you move on}]
What does \hi{सुनसान} mean? (\textquotedblleft{}Deserted\textquotedblright{}.) Say it once more.
\end{tcolorbox}
\section[\texorpdfstring{\hi{भीड़भाड़} (bhīṛbhāṛ)}{भीड़भाड़ (bhīṛbhāṛ)}]{\hi{भीड़भाड़} (bhīṛbhāṛ) — crowded, bustle}

\label{lesson:HI-C239-bhirbhar}



\begin{quote}
Before the new one: say the Hindi for hazy, then the Hindi for deserted.
\end{quote}

\subsection*{You'll want to know: \hi{भीड़भाड़}}
\textbf{\hi{भीड़भाड़}} — \emph{bhīṛbhāṛ} — \textquotedblleft{}crowded, bustle\textquotedblright{}.

\begin{cousinweb}
\textbf{\hi{भीड़भाड़}} (\emph{bhīṛbhāṛ}) is \emph{bhīṛ}, \textquotedblleft{}crowd\textquotedblright{}, which CDIAL lists under a reconstructed \emph{*bhiṭ}, \textquotedblleft{}meet, throng\textquotedblright{} (CDIAL 9490), with a rhyming partner \emph{bhāṛ}.
\end{cousinweb}

\begin{grammarlens}[title={What you've built}]
One more word for this chapter.
\end{grammarlens}

\subsection*{Your turn}
\begin{itemize}
\raggedright
  \item \emph{Say it:} \emph{bhīṛbhāṛ}
  \item \emph{Say it:} \emph{bhīṛbhāṛ}, once more
  \item \emph{Say it:} say \emph{bhīṛbhāṛ}
  \item \emph{Recall:} say the Hindi for hazy, then the Hindi for deserted, then say \emph{bhīṛbhāṛ} again
\end{itemize}

\begin{tcolorbox}[breakable,colback=teal!4,colframe=teal!35!black,title={Before you move on}]
What does \hi{भीड़भाड़} mean? (\textquotedblleft{}Crowded, bustle\textquotedblright{}.) Say it once more.
\end{tcolorbox}
\section[\texorpdfstring{\hi{आबाद} (ābād)}{आबाद (ābād)}]{\hi{आबाद} (ābād) — inhabited, lively}

\label{lesson:HI-C239-abad}



\begin{quote}
Before the new one: say the Hindi for deserted, then the Hindi for crowded.
\end{quote}

\subsection*{You'll want to know: \hi{आबाद}}
\textbf{\hi{आबाद}} — \emph{ābād} — \textquotedblleft{}inhabited, lively\textquotedblright{}.

\begin{cousinweb}
\textbf{\hi{आबाद}} (\emph{ābād}) is a loan, usually traced to Persian \emph{ābād}, \textquotedblleft{}inhabited, flourishing\textquotedblright{}.
\end{cousinweb}

\begin{grammarlens}[title={What you've built}]
One more word for this chapter.
\end{grammarlens}

\subsection*{Your turn}
\begin{itemize}
\raggedright
  \item \emph{Say it:} \emph{ābād}
  \item \emph{Say it:} \emph{ābād}, once more
  \item \emph{Say it:} say \emph{ābād}
  \item \emph{Recall:} say the Hindi for deserted, then the Hindi for crowded, then say \emph{ābād} again
\end{itemize}

\begin{tcolorbox}[breakable,colback=teal!4,colframe=teal!35!black,title={Before you move on}]
What does \hi{आबाद} mean? (\textquotedblleft{}Inhabited, lively\textquotedblright{}.) Say it once more.
\end{tcolorbox}
